\chapter{实数完备性的等价形式}
\label{chap:equivalent-completeness}

确界原理、单调有界收敛定理、区间套定理、Bolzano--Weierstrass 定理和 Cauchy 收敛准则的表述差别很大：它们分别谈论集合、单调数列、闭区间、子列和尾部误差。前几章已经逐一建立这些结论。本章不再重复全部证明，而要回答两个问题：这些结论怎样相互推出；在把其中一种形式当作公理时，还需要哪些序结构或 Archimedes 条件。

这种比较并非只为列出一组等价命题。在实际证明中，已知条件往往只适合其中一种形式。单调逼近适合取确界，逐次排除适合构造区间套，任意有界序列适合抽取子列，而误差估计通常直接给出 Cauchy 性。清楚各条推导所用的条件，才能在后文选择合适的完备性工具而不造成循环引用。

\section{五种完备性命题}

设 \(F\) 为有序域。绝对值、数列收敛、子列和 Cauchy 列均按前几章的定义理解。考虑以下命题：
\begin{description}[style=nextline,leftmargin=3em]
 \item[(S) 确界原理] 每个非空且有上界的集合 \(A\subseteq F\) 都有上确界。
 \item[(M) 单调有界收敛] 每个单调递增且有上界的 \(F\)-值数列都收敛；递减情形对称。
 \item[(N) 区间套] 若 \(I_n=[a_n,b_n]\) 满足 \(I_{n+1}\subseteq I_n\) 且 \(b_n-a_n\to0\)，则 \(\bigcap_n I_n\) 恰含一个点。
 \item[(B) Bolzano--Weierstrass] 每个有界 \(F\)-值数列都有收敛子列。
 \item[(C) Cauchy 完备] 每个 \(F\)-值 Cauchy 列都收敛。
\end{description}

另记
\begin{description}[style=nextline,leftmargin=3em]
 \item[(A) Archimedes 性质] 对每个 \(x\in F\)，存在 \(n\in\Npos\) 使 \(n>x\)。
\end{description}

\begin{theorem}[完备性命题的等价性]
在 Archimedes 有序域中，(S)、(M)、(N)、(B)、(C) 等价。确界原理 (S) 本身还蕴含 Archimedes 性质。
\end{theorem}

证明将在本章余下各节完成。可以采用许多不同的推导路线；这里选择
\[
 \mathrm{(S)}\Longrightarrow\mathrm{(M)}
 \Longrightarrow\mathrm{(N)}
 \Longrightarrow\mathrm{(B)}
 \Longrightarrow\mathrm{(C)},
 \qquad
 \mathrm{(A)}+\mathrm{(C)}\Longrightarrow\mathrm{(S)}.
\]
第一条链把集合的序边界逐步转化为数列的极限，最后一个方向则用二分法把任意集合的上边界编码成两条 Cauchy 列。长度 \(2^{-n}\) 能够小于任意正数，正是 Archimedes 性质介入的地方。

\section{从确界原理到 Cauchy 完备性}

\subsection{确界与单调数列}

设 \((a_n)\) 单调递增且有上界。由 (S)，值集
\[
 E=\{a_n:n\in\Npos\}
\]
有上确界 \(s\)。对任意 \(\varepsilon>0\)，\(s-\varepsilon\) 不是 \(E\) 的上界，故存在 \(N\) 使 \(a_N>s-\varepsilon\)。单调性把这一项的估计传给整个尾部：
\[
 s-\varepsilon<a_N\le a_n\le s\qquad(n\ge N).
\]
因此 \(a_n\to s\)。这正是\chapterref{chap:monotone-iteration}中单调有界收敛定理的证明；递减情形可对 \((-a_n)\) 应用同一结论。于是 (S) 推出 (M)。

这里两项假设承担不同工作。有界性使有限确界存在，单调性则保证进入 \(s\) 的一个邻域后不再离开。数列 \(a_n=n\) 说明不能删除有界性，\((-1)^n\) 说明不能删除单调性。

\subsection{单调收敛与区间套}

假设 (M) 成立，并设闭区间 \(I_n=[a_n,b_n]\) 嵌套且 \(b_n-a_n\to0\)。嵌套性给出
\[
 a_1\le a_2\le\cdots\le b_2\le b_1,
\]
所以 \((a_n)\) 单调递增且有上界，\((b_n)\) 单调递减且有下界。记
\[
 a_n\to x,\qquad b_n\to y.
\]
由极限的线性运算，\(b_n-a_n\to y-x\)，而已知左端趋于 \(0\)，故 \(x=y\)。对固定 \(n\)，当 \(k\ge n\) 时有
\[
 a_n\le a_k\le b_k\le b_n.
\]
令 \(k\to\infty\)，得到 \(a_n\le x\le b_n\)，因而 \(x\in I_n\) 对每个 \(n\) 成立。若 \(z\) 也属于全部 \(I_n\)，则
\[
 |z-x|\le b_n-a_n\qquad(n\in\Npos),
\]
长度趋零迫使 \(z=x\)。于是 (M) 推出 (N)。

长度条件只负责唯一性，并非闭区间套非空所必需。

\begin{proposition}[一般闭区间套]
设 \(I_n=[a_n,b_n]\) 为嵌套的非空闭区间。若 (S) 成立，则
\[
 \bigcap_{n=1}^{\infty}I_n
 =\left[\sup_n a_n,\inf_n b_n\right].
\]
特别地，交集非空，但可能含有多个点。
\end{proposition}

\begin{proof}
嵌套性使 \(a_n\le b_m\) 对任意 \(m,n\) 成立。故
\[
 a:=\sup_n a_n\le b:=\inf_n b_n.
\]
点 \(x\) 属于全部 \(I_n\)，当且仅当对每个 \(n\) 有 \(a_n\le x\le b_n\)，这又等价于 \(a\le x\le b\)。
\end{proof}

闭性不能省略，例如开区间 \((0,1/n)\) 的交为空；长度趋零不能保证非闭区间的端点被保留下来。

\subsection{区间套与收敛子列}

假设 (N) 成立，设有界数列 \((x_n)\) 的所有项落在闭区间 \(I_0=[A,B]\) 中。将 \(I_0\) 二分，至少有一个半区间含有无穷多个序列项；选定这样的半区间为 \(I_1\)。继续二分，得到嵌套闭区间
\[
 I_k=[\alpha_k,\beta_k],\qquad
 \beta_k-\alpha_k=2^{-k}(B-A),
\]
每个 \(I_k\) 都含有无穷多个序列项。Archimedes 性质保证区间长度趋于零，(N) 给出唯一交点 \(x\)。

递归选择 \(n_k>n_{k-1}\) 且 \(x_{n_k}\in I_k\)。这种选择可行，因为 \(I_k\) 中出现的指标有无穷多个。于是
\[
 |x_{n_k}-x|\le\beta_k-\alpha_k\longrightarrow0,
\]
从而得到收敛子列。故 (N) 推出 (B)。这也是\chapterref{chap:subsequences-cluster-points}中 Bolzano--Weierstrass 定理的区间套证明。

“含有无穷多个项”是关于指标的陈述，不是关于不同取值的陈述。常值数列的值集只有一个元素，却仍可在每一步抽取更大的指标。

\subsection{收敛子列与 Cauchy 收敛}

假设 (B) 成立，设 \((a_n)\) 为 Cauchy 列。它有界，因而存在子列 \(a_{n_k}\to a\)。给定 \(\varepsilon>0\)，先取 \(N\) 使
\[
 |a_m-a_n|<\frac{\varepsilon}{2}\qquad(m,n\ge N),
\]
再取 \(k\) 使 \(n_k\ge N\) 且 \(|a_{n_k}-a|<\varepsilon/2\)。于是对所有 \(n\ge N\)，
\[
 |a_n-a|
 \le |a_n-a_{n_k}|+|a_{n_k}-a|
 <\varepsilon.
\]
所以原列收敛到 \(a\)，即 (B) 推出 (C)。这一步把\chapterref{chap:cauchy-criterion}中的两项事实合在一起使用：Cauchy 列有界；Cauchy 列若有收敛子列，则全列收敛到同一点。

\section{从 Cauchy 完备性重建上确界}

前一节从 (S) 出发，逐步得到 (C)。反向推导不能简单地把箭头倒过来，因为 Cauchy 完备性只谈可数序列，而确界原理量化任意非空有界集合。需要用二分区间把集合的边界转化为可数数据。

\begin{theorem}
设 \(F\) 是 Archimedes 有序域。若 \(F\) 中每个 Cauchy 列都收敛，则 \(F\) 满足确界原理。
\end{theorem}

\begin{proof}
设 \(\varnothing\ne A\subseteq F\) 有上界。取 \(a_0\in A\)，令 \(l_0=a_0-1\)，则 \(l_0\) 不是上界；再取 \(A\) 的一个上界 \(u_0\)，必要时放大它，使 \(l_0<u_0\)。

已经构造 \(l_n,u_n\) 后，令 \(m_n=(l_n+u_n)/2\)。若 \(m_n\) 是 \(A\) 的上界，置
\[
 l_{n+1}=l_n,\qquad u_{n+1}=m_n;
\]
否则置
\[
 l_{n+1}=m_n,\qquad u_{n+1}=u_n.
\]
对每个 \(n\)，\(l_n\) 不是上界，\(u_n\) 是上界，并且
\[
 u_n-l_n=2^{-n}(u_0-l_0).
\]

Archimedes 性质使右端趋于零。若 \(m\ge n\)，嵌套性给出
\[
 l_n\le l_m\le u_m\le u_n,
\]
从而
\[
 |l_m-l_n|\le u_n-l_n,
 \qquad
 |u_m-u_n|\le u_n-l_n.
\]
所以 \((l_n)\)、\((u_n)\) 都是 Cauchy 列。由假设，存在 \(l,u\in F\) 使 \(l_n\to l\)、\(u_n\to u\)。又因 \(u_n-l_n\to0\)，得到 \(u=l=:s\)。

若有 \(a\in A\) 满足 \(a>s\)，取 \(\varepsilon=(a-s)/2\)。充分大时 \(u_n<s+\varepsilon<a\)，这与 \(u_n\) 是上界矛盾，故 \(s\) 是上界。若 \(c<s\)，充分大时 \(l_n>c\)；由于 \(l_n\) 不是上界，存在 \(a\in A\) 使 \(a>l_n>c\)，所以 \(c\) 不是上界。于是 \(s=\sup A\)。
\end{proof}

\begin{proposition}
在 Archimedes 有序域中，若每个单调递增有界数列都收敛，则确界原理成立。
\end{proposition}

\begin{proof}
对非空有上界集合 \(A\) 作上述二分构造。左端点列 \((l_n)\) 单调递增且有上界，故收敛到某个 \(s\)。区间宽度趋于零，因而 \(u_n\to s\)。再用“\(l_n\) 不是上界、\(u_n\) 是上界”两项不变量，重复上一证明的最后一段，即得 \(s=\sup A\)。
\end{proof}

\section{Archimedes 条件与有理数反例}

确界原理在有序域中能够推出 Archimedes 性质，证明已见\chapterref{chap:real-number-system}。反向从序列形式重建确界时，Archimedes 性质不能隐去。二分法所依赖的事实是
\[
 2^{-n}(u_0-l_0)\longrightarrow0.
\]
在非 Archimedes 有序域中，可能存在正无穷小 \(\eta\)，使 \(2^{-n}>\eta\) 对每个自然数 \(n\) 都成立；此时可数次二分未必越过所有尺度。

\begin{proposition}[同一个有理缺口的五种表现]
\(\Q\) 是 Archimedes 有序域，但 (S)、(M)、(N)、(B)、(C) 全部失败。
\end{proposition}

\begin{proof}
令
\[
 E=\{q\in\Q:q<0\text{ 或 }q^2<2\}.
\]
由\chapterref{chap:real-number-system}，\(E\) 在 \(\Q\) 中有上界却没有有理上确界，故 (S) 失败。取从左侧逼近 \(\sqrt2\) 的十进制截断列 \((a_n)\)，它单调递增且有界，却没有有理极限，故 (M) 失败。相应的左右有理端点构成长趋于零的闭区间套，其交在 \(\Q\) 中为空，故 (N) 失败。

同一列是有理 Cauchy 列而无有理极限，所以 (C) 失败。若 (B) 成立，\((a_n)\) 应有在 \(\Q\) 中收敛的子列；设该子列趋于 \(q\in\Q\)，则平方的极限给出 \(q^2=2\)，仍与 \(\sqrt2\notin\Q\) 矛盾。故 (B) 也失败。
\end{proof}

这个例子说明 Archimedes 性质只保证自然数网格可以任意细分，并不保证网格逼近的边界仍属于原数系。完备性补充的是极限留在空间内的封闭性。

\begin{remark}[序完备与 Cauchy 完备]
确界原理检查任意非空有界子集的序边界，Cauchy 完备性只检查可数序列的距离稳定性。二者在 Archimedes 有序域中由上面的二分编码联系起来；在一般度量空间中没有上确界或区间可用，Cauchy 完备性仍有意义，而本章的其他形式必须另作解释或替换。
\end{remark}

\section{在证明中选择完备性形式}

等价性不意味着各形式在证明中同样方便。若对象天然单调，例如积分上和、下和或逐步增大的近似集合，通常直接寻找上确界或下确界。若每一步把候选范围缩小到闭区间，区间套定理同时记录存在性和定位误差。若只有有界性而缺少单调性，Bolzano--Weierstrass 定理允许先抽取收敛子列；若能估计任意两个后期近似之间的差，Cauchy 准则往往无需预知极限。

还要区分“存在候选值”和“候选值属于当前空间”。极限方程、上下极限或形式代数运算可以给出候选值，却不能代替完备性论证。反过来，完备性通常只给出存在性；极限的具体数值还要靠唯一性、方程、保序性或误差估计确定。后续章节使用完备性时，将明确指出采用的是上述哪一种形式。

\section*{习题}
\addcontentsline{toc}{section}{习题}

\subsection*{A组　定义与基本方法}
\begin{enumerate}[label=\textbf{\thechapter.\arabic*.}]
 \exerciseitem{6.1} 从定义证明收敛数列是 Cauchy 列。
 \exerciseitem{6.2} 证明单调递增列若无上界，则对每个 \(M\in\R\) 从某项起 \(a_n>M\)。
 \exerciseitem{6.3} 给出有界但不收敛数列、单调但不收敛数列，并指出各自缺少的条件。
 \exerciseitem{6.4} 证明区间套定理中若只知交集非空而无长度条件，交集必为闭区间。
 \exerciseitem{6.5} 分别给出删除嵌套性、闭性、长度趋零条件后结论失败的例子。
\end{enumerate}

\subsection*{B组　证明与反例}
\begin{enumerate}[label=\textbf{\thechapter.\arabic*.},resume]
 \exerciseitem{6.6} 证明有界数列若只有一个可能的子列极限，则原列收敛。
 \exerciseitem{6.7} 证明 Cauchy 列的任意子列仍是 Cauchy 列；若某子列收敛，则所有收敛子列极限相同。
 \exerciseitem{6.8} 用二分区间法证明每个无限有界实数集至少有一个聚点。
 \exerciseitem{6.9} 证明若 \(|a_{n+1}-a_n|\le2^{-n}\)，则 \((a_n)\) 是 Cauchy 列。
 \optionalexerciseitem{6.10} 构造 \(|a_{n+1}-a_n|\to0\) 但不是 Cauchy 列的数列。
 \exerciseitem{6.11} 证明有限多个 Cauchy 列的线性组合仍为 Cauchy 列。
\end{enumerate}

\subsection*{C组　综合问题}
\begin{enumerate}[label=\textbf{\thechapter.\arabic*.},resume]
 \projectexerciseitem{6.12} \ 【前置：本章全部主定理；预计用时：90分钟】不经过 Bolzano--Weierstrass 定理，直接用快速 Cauchy 子列与级数型误差证明 Cauchy 完备性。
 \exerciseitem{6.13} 补全从 Cauchy 完备性构造上确界时 \(l_n,u_n\) 的 Cauchy 性和共同极限证明。
 \exerciseitem{6.14} 证明“每个有界数列都有 Cauchy 子列”与 Bolzano--Weierstrass 定理在 Cauchy 完备的有序域中等价。
 \exerciseitem{6.15} 设闭区间套 \(I_n=[a_n,b_n]\) 的长度趋于零，交点为 \(x\)。证明对任意选择 \(x_n\in I_n\)，都有 \(x_n\to x\)。
 \optionalexerciseitem{6.16} 设有序域 \(F\) 中存在正元素 \(\eta\)，满足 \(\eta<1/n\) 对每个 \(n\in\Npos\) 成立。证明 \(2^{-n}\) 在 \(F\) 中不趋于零，并说明可数次二分为何不能越过尺度 \(\eta\)。
 \exerciseitem{6.17} 证明在 \(\Q\) 中 (S)、(M)、(N)、(B)、(C) 均失败，并尽量使用\chapterref{chap:real-number-system}同一个 \(\sqrt2\) 缺口。
 \exerciseitem{6.18} 对不要求长度趋零的嵌套闭区间证明
 \[
  \bigcap_n[a_n,b_n]=[\sup_n a_n,\inf_n b_n].
 \]
 \exerciseitem{6.19} 从 Bolzano--Weierstrass 定理推出每个无限有界实数集有聚点，并解释互异点选择的必要性。
 \exerciseitem{6.20} 定义尾部直径 \(d_N\)，证明 Cauchy 性等价于 \(d_N\to0\)。
 \exerciseitem{6.21} 证明若 \((a_n)\) 是 Cauchy 列且有子列趋于 \(+\infty\)，则产生矛盾；说明 Cauchy 列不能有不同的扩充实子列极限。
 \exerciseitem{6.22} 只用单调有界收敛与 Archimedes 性质重建确界原理。
 \exerciseitem{6.23} 证明若每个长度趋零的闭区间套交集非空，则交集自动为单点；据此区分“存在性”与“唯一性”的条件来源。
 \optionalexerciseitem{6.24} 给出从 (C) 直接推出 (B) 的证明方案，避免先使用确界原理；可先构造 Cauchy 子列。
 \projectexerciseitem{6.25} \ 【前置：本章全部内容；预计用时：120分钟】为依赖图中每条箭头写一份条件清单和反例清单，确保不循环引用等价命题。
 \optionalexerciseitem{6.26} 设 \(A\subseteq\R\) 非空且有上界，\(s\in\R\)。证明：\(s=\sup A\) 当且仅当 \(s\) 是 \(A\) 的上界，并且存在 \(a_n\in A\) 使 \(a_n\to s\)。
\end{enumerate}

\section*{部分习题提示}
\begin{itemize}
 \item \exerciseref{6.6}：反设不收敛，抽取一条始终远离候选极限的子列，再应用收敛子列定理。
 \item \exerciseref{6.10}：取调和部分和 \(a_n=\sum_{k=1}^n1/k\)。
 \item \exerciseref{6.12}：抽取 \(n_k\) 使尾部项两两距离 \(<2^{-k}\)，证明它收敛后再控制原列。
\end{itemize}
